\section{Related Work}
\label{sec:related}

\paragraph{Temporal shifting.}
Wiesner et al.~\cite{wiesner2021wait} showed that deferring
delay-tolerant jobs to low-carbon hours reduces emissions, with gains
that depend on the region and the length of the flexibility window.
Google shapes datacenter capacity using day-ahead carbon
forecasts~\cite{radovanovic2023carbon}, and CarbonScaler~\cite{hanafy2023carbonscaler}
scales elastic jobs with carbon intensity. These works take a single
region as given. GreenOps uses Wiesner et al.'s strategy as a baseline.

\paragraph{Spatial and spatio-temporal shifting.}
Geo-distributed carbon-aware load balancing dates back at least to
Zhou et al.~\cite{zhou2013carbon}, and has since been applied to
serverless functions~\cite{chadha2023greencourier,gsteiger2024caribou}
and edge computing~\cite{wu2025carbonedge}. Sukprasert et
al.~\cite{sukprasert2024spatiotemporal} extend joint shifting to water
and land use. Lechowicz et al.~\cite{lechowicz2024soad} give
competitive online algorithms for spatio-temporal allocation with
deadlines and movement costs. These build on threshold
rules~\cite{elyaniv2001search}, which we use as a forecast-free baseline.

\paragraph{Limits of shifting.}
Acun et al.~\cite{acun2024limitations} found that the realistic
benefits of both kinds of shifting are much smaller than idealised
analyses suggest. Simple policies capture most of what is achievable,
and capacity, latency and the cost of moving work limit the rest. Our
study is complementary: we measure how SLA risk, data movement and
forecast error shape the trade-off in a trace-driven setting, and
compare published policies under identical constraints.
